\subsection{模型优点}

本文模型的优势在于把题目的三层要求分别交给相应的数学工具处理：
\begin{enumerate}[label=(\arabic*),leftmargin=2.5em]
\item 缺失感知评价避免将结构性缺失或未公开成绩误判为低能力；
\item Benchmark Family 平衡抑制测试数量造成的隐性重复赋权；
\item Bradley--Terry 成对比较将不同量纲的公开测试转化为统一的潜在能力尺度；
\item CES 聚合描述任务中的能力互补和短板效应，避免线性加权的完全补偿；
\item 工作负载成本把 API 单价转化为实际任务成本；
\item Pareto 方法保留效用和成本的双重目标，避免人为压缩成一个缺乏解释力的分数；
\item Bootstrap、LOFO、LOSO 和参数敏感性从不同角度检验结论对数据与设定变化的响应。
\end{enumerate}

\subsection{局限与改进方向}

第一，公开 Benchmark 是横截面信息，不能完全替代统一实验环境，模型版本和测试协议变化也可能影响横向比较。第二，Bradley--Terry 能力估计依赖可观测比较网络；当桥接设置缺失时，部分维度可能出现网络断开或结构依赖。第三，C5 的有效模型范围较小，相关结论应与适用性审计一并解释。第四，Q2 不确定性传播采用正式接口给出的独立正态近似场景样本，而不是带完整协方差的联合 Q1 Bootstrap。第五，Q3 使用标准化工作负载，不能替代特定用户的真实 Token 分布；价格也可能随时间变化。最后，成本信息不完整的模型无法进入严格 Pareto，这体现了部署决策的可核验性要求，也限制了效用与成本的直接覆盖范围。

后续可在统一评测协议下补充联合 Bootstrap 协方差、时间序列价格、真实 Token 负载、延迟和吞吐等部署指标，并在保留 Pareto 结构的前提下扩展多目标决策。
